\documentclass[11pt,a4paper]{article}

\usepackage[utf8]{inputenc}
\usepackage[T1]{fontenc}
\usepackage[portuguese]{babel}
\usepackage[a4paper,margin=2.25cm,headheight=14pt]{geometry}
\usepackage[table]{xcolor}
\usepackage{booktabs}
\usepackage{tabularx}
\usepackage{array}
\usepackage{enumitem}
\usepackage{fancyhdr}
\usepackage{hyperref}

\definecolor{petrolblue}{RGB}{15,50,90}
\definecolor{petrolgold}{RGB}{200,160,50}
\definecolor{lightgray}{RGB}{245,247,250}

\hypersetup{
  colorlinks=true,
  linkcolor=petrolblue,
  urlcolor=petrolblue,
  pdftitle={PetroLúmen — Modelo de Negócio GEOENERGY 2026},
  pdfsubject={Anexo I — Modelo de Negócio},
  pdfauthor={Equipa PetroLúmen}
}

\pagestyle{fancy}
\fancyhf{}
\fancyhead[L]{\small\textcolor{petrolblue}{PetroLúmen}}
\fancyhead[R]{\small\textcolor{petrolblue}{GEOENERGY 2026}}
\fancyfoot[C]{\thepage}
\renewcommand{\headrulewidth}{0.4pt}

\setlength{\parindent}{0pt}
\setlength{\parskip}{0.45em}
\setlength{\emergencystretch}{3em}
\renewcommand{\arraystretch}{1.18}
\setlist[itemize]{leftmargin=*,topsep=2pt,itemsep=2pt}
\setlist[enumerate]{leftmargin=*,topsep=2pt,itemsep=3pt}

\newcommand{\usd}{US\$}
\newcommand{\confirmar}{\textit{(a confirmar pela equipa)}}

\begin{document}

\pagenumbering{gobble}
\begin{titlepage}
  \thispagestyle{empty}
  \centering
  \vspace*{1.7cm}
  {\Large\textcolor{petrolgold}{GEOENERGY — 2.ª EDIÇÃO 2026}\par}
  \vspace{1.5cm}
  {\Huge\bfseries\textcolor{petrolblue}{PETROLÚMEN}\par}
  \vspace{0.5cm}
  {\Large Plataforma de Engenharia de Reservatórios\par}
  \vspace{1.5cm}
  \rule{\textwidth}{1.2pt}\par
  \vspace{0.4cm}
  {\LARGE\bfseries Modelo de Negócio\par}
  \vspace{0.4cm}
  \rule{\textwidth}{1.2pt}\par
  \vspace{1.2cm}
  {\large Anexo I — estrutura recomendada no Regulamento GEOENERGY\par}
  \vfill
  \begin{tabularx}{0.82\textwidth}{>{\bfseries}l X}
    Equipa & Rocélio Da Silva; Derik Francisco; Davidaldo Belchior; Mafra Tito \\
    Proponente principal & Rocélio Da Silva \\
    Instituição indicada no material da candidatura & ISPTEC \\
    Categoria proposta & Digitalização e Tecnologias Aplicadas às Geociências e Energia \\
    Data & Outubro de 2026 \\
  \end{tabularx}
  \vfill
  \begin{minipage}{0.86\textwidth}
    \small
    \textit{Documento de candidatura. Os dados de elegibilidade da equipa, a procura,
    os preços, os custos e as projeções são hipóteses preliminares e devem ser
    confirmados pelos proponentes antes da submissão.}
  \end{minipage}
  \vspace{0.8cm}
\end{titlepage}
\clearpage
\pagenumbering{arabic}

\section{Identificação da equipa e do proponente principal}

\begin{tabularx}{\textwidth}{>{\bfseries}p{4.1cm} X}
  Proponente principal & Rocélio Da Silva \\
  Membros da equipa & Derik Francisco; Davidaldo Belchior; Mafra Tito \\
  Instituição & ISPTEC \\
  Curso declarado no rascunho & Engenharia de Petróleos \\
  Contactos e números de matrícula & A indicar na ficha de inscrição \\
\end{tabularx}

O material de candidatura identifica os proponentes como estudantes de Engenharia de
Petróleos do ISPTEC. A equipa deve confirmar a situação académica, os números de
matrícula e que o proponente principal satisfaz os requisitos de elegibilidade do
regulamento. A candidatura deve ser subscrita por todos os membros.

\section{Título da ideia de negócio}

\textbf{PetroLúmen — Plataforma Integrada de Engenharia de Reservatórios.}

O nome é curto e diretamente associado à energia e à engenharia. A equipa deve
verificar a disponibilidade do nome comercial e dos identificadores digitais antes
de iniciar a exploração comercial.

\section{Problema ou oportunidade identificada}

Estudantes, engenheiros e consultores de reservatórios recorrem frequentemente a
folhas de cálculo, scripts independentes e ferramentas que não partilham dados nem
fluxos de trabalho. Esta fragmentação torna análises repetitivas, aumenta o esforço
de verificação e dificulta a utilização consistente de métodos de engenharia.

As soluções comerciais especializadas podem estar fora do alcance de estudantes,
consultores independentes e pequenas organizações. Existe, por isso, uma oportunidade
para uma ferramenta de apoio à análise técnica que seja acessível, modular e adaptada
a fluxos de trabalho de engenharia de reservatórios.

Esta oportunidade é uma hipótese de negócio, não uma medição de mercado. Deve ser
validada por entrevistas a utilizadores, docentes, consultores e potenciais clientes
em Angola.

\section{Solução proposta e proposta de valor}

O PetroLúmen é uma aplicação de engenharia de reservatórios que reúne ferramentas de
análise num ambiente de trabalho comum. A base de código existente combina uma
interface em React/Next.js, empacotamento desktop com Tauri e um backend Python/FastAPI.
O repositório inclui módulos e endpoints para cálculos PVT, análise de produção e
declínio, análise nodal/RTA, simulação de reservatórios, previsão e gestão de dados.

\textbf{Valor proposto:} reduzir a fragmentação do trabalho analítico, tornar
ferramentas de engenharia mais acessíveis e facilitar a repetibilidade dos fluxos
de análise. O software destina-se a apoiar o trabalho técnico; os resultados devem
ser revistos por profissionais qualificados e não substituem validação de engenharia.

\textbf{Diferencial pretendido:} foco em tarefas frequentes de reservatórios, interface
integrada, evolução modular e possibilidade de implantação local. A operação
completamente offline, a segurança de dados e a integração de todos os módulos devem
ser demonstradas em testes e pilotos antes de serem anunciadas como garantias.

\section{Mercado e clientes-alvo}

\begin{itemize}
  \item \textbf{Estudantes e instituições de ensino:} uso pedagógico, exercícios e
  demonstrações de análise de reservatórios.
  \item \textbf{Engenheiros e consultores independentes:} apoio a análises de rotina,
  preparação de cenários e elaboração de resultados técnicos.
  \item \textbf{Pequenas consultoras e operadores:} ferramenta modular para equipas
  que necessitam de fluxos de análise acessíveis e configuráveis.
  \item \textbf{Empresas de maior dimensão:} possibilidade futura de implantação
  local, sujeita a requisitos de segurança, integração, suporte e validação.
\end{itemize}

A entrada no mercado deve começar por entrevistas e demonstrações junto da comunidade
académica e de profissionais locais. O número de utilizadores potencialmente
interessados, a disposição para pagar e os requisitos de aquisição empresarial ainda
precisam de ser medidos. Não se assume qualquer contrato ou parceria existente.

\section{Diferencial competitivo e análise sumária da concorrência}

\begin{tabularx}{\textwidth}{>{\bfseries}p{3.2cm} X X}
  \toprule
  Alternativa & Vantagem percebida & Espaço de diferenciação do PetroLúmen \\
  \midrule
  Folhas de cálculo e scripts próprios
    & Flexíveis e já conhecidos pelos utilizadores
    & Reunir fluxos de análise, facilitar reutilização e reduzir tarefas manuais \\
  Ferramentas científicas de código aberto
    & Acesso ao código e possibilidade de adaptação
    & Oferecer interface de produto e fluxos mais orientados ao utilizador \\
  Suites comerciais especializadas
    & Funcionalidades maduras e suporte estabelecido
    & Explorar acessibilidade, modularidade e necessidades do contexto local \\
  \bottomrule
\end{tabularx}

Esta comparação é qualitativa e não substitui uma análise de preços, licenças,
funcionalidades e desempenho de concorrentes específicos. A diferenciação deverá ser
confirmada com utilizadores e não deve assentar apenas em alegações de menor custo ou
maior precisão.

\section{Modelo de receitas ou mecanismo de criação de valor}

O modelo de receitas é uma hipótese escalonada, ainda por validar:

\begin{enumerate}
  \item \textbf{Plano académico:} versão gratuita ou de baixo custo para estudantes,
  com limites de utilização ou exportação a definir.
  \item \textbf{Plano profissional:} subscrição para consultores e profissionais,
  com acesso a funcionalidades avançadas e suporte.
  \item \textbf{Plano organizacional:} licença ou contrato anual para organizações,
  potencialmente com implantação local, configuração e formação.
\end{enumerate}

O rascunho comercial considera, para simulação, \usd 40 por mês no plano profissional
e \usd 4\,000 por ano no plano organizacional. Estes valores são referências
hipotéticas, não preços aprovados nem validados por clientes. O mecanismo de
licenciamento, faturação, suporte e conformidade deverá ser definido antes do
lançamento.

\section{Viabilidade técnica e recursos necessários}

\textbf{Base existente no repositório:}
\begin{itemize}
  \item frontend em React/Next.js e configuração de aplicação desktop Tauri;
  \item backend FastAPI organizado em API e módulos Python de engenharia;
  \item funcionalidades de cálculo PVT, análise/declínio de produção, previsão,
  simulação e endpoints relacionados;
  \item testes automatizados organizados em \texttt{backend/tests/} e testes de
  interface configurados com Vitest.
\end{itemize}

A existência de código e testes não equivale, por si só, a certificação científica,
prontidão industrial ou validação independente dos cálculos. Antes de um piloto,
deverão ser verificados os resultados contra exemplos de referência, a integração
entre a aplicação desktop e o backend, o tratamento de dados, a documentação de
limitações e a experiência de utilização.

\textbf{Recursos necessários:} tempo de desenvolvimento da equipa; revisão técnica
por docentes e engenheiros de reservatórios; equipamentos de desenvolvimento e
demonstração; testes com dados sintéticos ou autorizados; apoio de UX, validação de
segurança e orientação comercial. Dados reais de poços só devem ser utilizados com
autorização e controlos adequados de confidencialidade.

\section{Estimativa económico-financeira preliminar}

A tabela seguinte é um cenário de referência, não uma previsão garantida. Assume
subscrições profissionais ao preço hipotético de \usd 40/mês durante 12 meses e
licenças organizacionais de \usd 4\,000/ano. Os custos são estimativas preliminares
do rascunho e carecem de cotações. Não incluem remuneração integral dos fundadores,
impostos, comissões, custos financeiros nem contingências.

\begin{table}[ht]
\centering
\small
\begin{tabularx}{\textwidth}{X r r r}
  \toprule
  \textbf{Indicador} & \textbf{Ano 1} & \textbf{Ano 2} & \textbf{Ano 3} \\
  \midrule
  Subscrições profissionais (contas-ano) & 50 & 150 & 350 \\
  Receita profissional (\usd) & 24\,000 & 72\,000 & 168\,000 \\
  Licenças organizacionais (unidades-ano) & 0 & 2 & 5 \\
  Receita organizacional (\usd) & 0 & 8\,000 & 20\,000 \\
  \textbf{Receita bruta modelada (\usd)} & \textbf{24\,000} & \textbf{80\,000} & \textbf{188\,000} \\
  Custos operacionais preliminares (\usd) & 5\,000 & 12\,000 & 25\,000 \\
  \textbf{Saldo operacional simplificado (\usd)} & \textbf{19\,000} & \textbf{68\,000} & \textbf{163\,000} \\
  \bottomrule
\end{tabularx}
\caption{Cenário ilustrativo baseado em hipóteses ainda não validadas.}
\end{table}

O saldo simplificado corresponde apenas à receita bruta menos os custos indicados;
não deve ser interpretado como lucro líquido. As quantidades de clientes, preços,
custos e calendário de aquisição devem ser revistos com dados de entrevistas, pilotos
e cotações. Não se afirma um mês de equilíbrio financeiro sem um fluxo de caixa
mensal validado.

\section{Impacto socioeconómico e ambiental}

\textbf{Socioeconómico:} o desenvolvimento local de software técnico pode fortalecer
competências digitais e de engenharia, criar propriedade intelectual e abrir
oportunidades de prestação de serviços e emprego qualificado. Estes impactos
dependem da adoção, continuidade do produto e capacidade de gerar receitas.

\textbf{Ambiental:} ferramentas de análise podem apoiar a comparação de cenários e
decisões de gestão de reservatórios. O software não reduz emissões ou perfurações
por si só; qualquer benefício ambiental deve ser medido num piloto, com indicadores
definidos e atribuição prudente dos resultados.

\textbf{Sustentabilidade e responsabilidade:} a proposta pretende criar valor
económico através de receitas recorrentes que financiem manutenção, validação e
suporte, e valor social através do desenvolvimento de competências locais. A
utilização deve respeitar a confidencialidade dos dados, a segurança dos utilizadores,
a responsabilidade profissional e a protecção ambiental. Estes resultados dependem
da implementação e serão acompanhados por indicadores de adopção, continuidade das
receitas, actividades de formação e evidência dos pilotos.

\section{Plano preliminar de implementação, riscos e mitigação}

\begin{tabularx}{\textwidth}{>{\bfseries}p{2.5cm} X X}
  \toprule
  Fase & Actividades e resultado esperado & Critério de avanço \\
  \midrule
  1. Validação
    & Entrevistar utilizadores, escolher tarefas prioritárias e confirmar requisitos
    & Problemas recorrentes e utilizadores-piloto identificados \\
  2. Consolidação técnica
    & Rever cálculos, integração, documentação, importação/exportação e estabilidade
    & Casos de referência reproduzidos e limitações documentadas \\
  3. Piloto académico
    & Demonstrar fluxos selecionados com dados autorizados ou sintéticos
    & Feedback recolhido, erros prioritários tratados \\
  4. Preparação comercial
    & Testar preços, suporte, licenciamento e custos de implantação
    & Interesse e disposição para pagar avaliados \\
  5. Lançamento gradual
    & Disponibilizar versão inicial e acompanhar utilização e suporte
    & Indicadores de uso, retenção e custos acompanhados \\
  Preparação para exposição
    & Preparar banner e demonstração segura dentro do espaço atribuído; levar os
    equipamentos necessários e solicitar autorização prévia para demonstrações
    & Cumprimento das instruções da organização e dos requisitos de segurança \\
  \bottomrule
\end{tabularx}

\vspace{0.5em}
\textbf{Principais riscos e mitigação:}
\begin{itemize}
  \item \textbf{Resultados técnicos incorretos:} testar contra exemplos de referência,
  rever unidades e hipóteses e apresentar limitações junto dos resultados.
  \item \textbf{Baixa adesão ou disposição para pagar:} validar necessidades e preços
  com entrevistas e pilotos antes de investir no lançamento.
  \item \textbf{Atrasos de integração e estabilidade:} reduzir o âmbito inicial a
  fluxos prioritários e manter testes de regressão.
  \item \textbf{Exposição de dados confidenciais:} começar com dados sintéticos ou
  autorizados, limitar acessos e validar armazenamento, transferência e eliminação.
  \item \textbf{Dependência da equipa e falta de recursos:} documentar o produto,
  priorizar entregas e procurar mentoria técnica e empresarial.
  \item \textbf{Custos e receitas diferentes das hipóteses:} obter cotações, rever
  periodicamente o orçamento e condicionar a expansão à validação comercial.
\end{itemize}

\section*{Declaração e subscrição da equipa}
\addcontentsline{toc}{section}{Declaração e subscrição da equipa}

Os proponentes abaixo identificados declaram que leram e aceitam o Regulamento
GEOENERGY 2026, que as informações fornecidas nesta candidatura são verdadeiras
segundo o seu conhecimento e que assumem conjuntamente a responsabilidade pela
proposta. Confirmam também que a equipa e a elegibilidade cumprem os requisitos
aplicáveis do regulamento. Cada membro deverá assinar antes da submissão.

\vspace{1em}
\begin{tabularx}{\textwidth}{X X}
  \textbf{Proponente principal} & \textbf{Membro da equipa} \\
  Nome: Rocélio Da Silva & Nome: Derik Francisco \\
  Assinatura: \hrulefill & Assinatura: \hrulefill \\
  Data: \hrulefill & Data: \hrulefill \\
\end{tabularx}

\vspace{1em}
\textit{Se houver outros membros, acrescentar os respectivos nomes, assinaturas e
datas. A candidatura deve ser subscrita por todos os elementos da equipa.}

\end{document}
